\documentclass{amsart}
% For dealing with references we use the comment environment
\usepackage{verbatim}
\newenvironment{reference}{\comment}{\endcomment}
%\newenvironment{reference}{}{}
\newenvironment{slogan}{\comment}{\endcomment}
\newenvironment{history}{\comment}{\endcomment}

% For commutative diagrams we use Xy-pic
\usepackage[all]{xy}

% We use 2cell for 2-commutative diagrams.
\xyoption{2cell}
\UseAllTwocells

% We use multicol for the list of chapters between chapters
\usepackage{multicol}

% This is generally recommended for better output
\usepackage{lmodern}
\usepackage[T1]{fontenc}

% For cross-file-references

% Package for hypertext links:
\usepackage{hyperref}

% For any local file, say "hello.tex" you want to link to please

% Theorem environments.
%
\theoremstyle{plain}
\newtheorem{theorem}[subsection]{Theorem}
\newtheorem{proposition}[subsection]{Proposition}
\newtheorem{lemma}[subsection]{Lemma}

\theoremstyle{definition}
\newtheorem{definition}[subsection]{Definition}
\newtheorem{example}[subsection]{Example}
\newtheorem{exercise}[subsection]{Exercise}
\newtheorem{situation}[subsection]{Situation}

\theoremstyle{remark}
\newtheorem{remark}[subsection]{Remark}
\newtheorem{remarks}[subsection]{Remarks}

\numberwithin{equation}{subsection}

% Macros
%
\def\lim{\mathop{\mathrm{lim}}\nolimits}
\def\colim{\mathop{\mathrm{colim}}\nolimits}
\def\Spec{\mathop{\mathrm{Spec}}}
\def\Hom{\mathop{\mathrm{Hom}}\nolimits}
\def\Ext{\mathop{\mathrm{Ext}}\nolimits}
\def\SheafHom{\mathop{\mathcal{H}\!\mathit{om}}\nolimits}
\def\SheafExt{\mathop{\mathcal{E}\!\mathit{xt}}\nolimits}
\def\Sch{\mathit{Sch}}
\def\Mor{\mathop{\mathrm{Mor}}\nolimits}
\def\Ob{\mathop{\mathrm{Ob}}\nolimits}
\def\Sh{\mathop{\mathit{Sh}}\nolimits}
\def\NL{\mathop{N\!L}\nolimits}
\def\CH{\mathop{\mathrm{CH}}\nolimits}
\def\proetale{{pro\text{-}\acute{e}tale}}
\def\etale{{\acute{e}tale}}
\def\QCoh{\mathit{QCoh}}
\def\Ker{\mathop{\mathrm{Ker}}}
\def\Im{\mathop{\mathrm{Im}}}
\def\Coker{\mathop{\mathrm{Coker}}}
\def\Coim{\mathop{\mathrm{Coim}}}

% Boxtimes
%
\DeclareMathSymbol{\boxtimes}{\mathbin}{AMSa}{"02}

%
% Macros for moduli stacks/spaces
%
\def\QCohstack{\mathcal{QC}\!\mathit{oh}}
\def\Cohstack{\mathcal{C}\!\mathit{oh}}
\def\Spacesstack{\mathcal{S}\!\mathit{paces}}
\def\Quotfunctor{\mathrm{Quot}}
\def\Hilbfunctor{\mathrm{Hilb}}
\def\Curvesstack{\mathcal{C}\!\mathit{urves}}
\def\Polarizedstack{\mathcal{P}\!\mathit{olarized}}
\def\Complexesstack{\mathcal{C}\!\mathit{omplexes}}
% \Pic is the operator that assigns to X its picard group, usage \Pic(X)
% \Picardstack_{X/B} denotes the Picard stack of X over B
% \Picardfunctor_{X/B} denotes the Picard functor of X over B
\def\Pic{\mathop{\mathrm{Pic}}\nolimits}
\def\Picardstack{\mathcal{P}\!\mathit{ic}}
\def\Picardfunctor{\mathrm{Pic}}
\def\Deformationcategory{\mathcal{D}\!\mathit{ef}}

\setlength{\emergencystretch}{3em}
\begin{document}
\title[Stacks Project, Tag 09AS]{355 --- Gluing perfectness across two quotient ideals}
\maketitle
\noindent Copyright (C) 2005--2025 Johan de Jong. Stacks Project authors. Source: \href{https://stacks.math.columbia.edu/tag/09AS}{Tag 09AS}.

\noindent Editorial difficulty note (not author text). Difficulty H2: After removing a square-zero quotient the proof realizes the object in a fibre-product exact sequence of K-flat complexes. It proves compactness through derived Hom comparisons over the three quotients and a five-lemma argument, recovering perfectness over the product ideal..

\noindent Editorial provenance note (not author text). History: excerpt from revision \texttt{a0444\allowbreak 6e57e\allowbreak c1fbc\allowbreak 252a8\allowbreak 71afc\allowbreak ec775\allowbreak 2fb28\allowbreak 07b14}, more-algebra.tex, lines 20531--20642. Reference presentation was adapted; original-author context and core support blocks were retained or added. The mathematical wording is unchanged. Licensed under GNU FDL 1.2 or later; no invariant sections or cover texts. See ../licenses/COPYING. Context blocks reproduce author definitions and supporting results; they are not additional counted problems.

\begin{lemma}
\label{lemma-perfect-modulo-two-ideals}
Let $R$ be a ring. Let $I, J \subset R$ be ideals.
Let $K$ be an object of $D(R)$. Assume that
\begin{enumerate}
\item $K \otimes_R^\mathbf{L} R/I$ is perfect in $D(R/I)$, and
\item $K \otimes_R^\mathbf{L} R/J$ is perfect in $D(R/J)$.
\end{enumerate}
Then $K \otimes_R^\mathbf{L} R/IJ$ is perfect in $D(R/IJ)$.
\end{lemma}

\begin{proof}
It is clear that we may assume replace $R$ by $R/IJ$ and $K$ by
$K \otimes_R^\mathbf{L} R/IJ$. Then $R \to R/(I \cap J)$ is
a surjection whose kernel has square zero. Hence by
Lemma \href{https://stacks.math.columbia.edu/tag/07LU}{lemma perfect modulo nilpotent ideal (Tag 07LU)}
it suffices to prove that $K \otimes_R^\mathbf{L} R/(I \cap J)$ is
perfect. Thus we may assume that $I \cap J = 0$.

\medskip\noindent
We prove the lemma in case $I \cap J = 0$. First, we may represent $K$
by a K-flat complex $K^\bullet$ with all $K^n$ flat, see
Lemma \href{https://stacks.math.columbia.edu/tag/06Y4}{lemma K flat resolution (Tag 06Y4)}. Then we see that we have a short
exact sequence of complexes
$$
0 \to
K^\bullet \to K^\bullet/IK^\bullet \oplus K^\bullet/JK^\bullet \to
K^\bullet/(I + J)K^\bullet \to 0
$$
Note that $K^\bullet/IK^\bullet$ represents $K \otimes^\mathbf{L}_R R/I$
by construction of the derived tensor product. Similarly for
$K^\bullet/JK^\bullet$ and $K^\bullet/(I + J)K^\bullet$.
Note that $K^\bullet/(I + J)K^\bullet$ is a perfect complex
of $R/(I + J)$-modules, see
Lemma \href{https://stacks.math.columbia.edu/tag/066W}{lemma pull perfect (Tag 066W)}.
Hence the complexes $K^\bullet/IK^\bullet$, and
$K^\bullet/JK^\bullet$ and $K^\bullet/(I + J)K^\bullet$
have finitely many nonzero cohomology groups
(since a perfect complex has finite Tor-amplitude, see
Lemma \href{https://stacks.math.columbia.edu/tag/0658}{lemma perfect (Tag 0658)}). We conclude that $K \in D^b(R)$ by the
long exact cohomology sequence associated to short exact sequence
of complexes displayed above. In particular we assume $K^\bullet$
is a bounded above complex of free $R$-modules (see
Derived Categories, Lemma \href{https://stacks.math.columbia.edu/tag/05T7}{lemma subcategory left resolution (Tag 05T7)}).

\medskip\noindent
We will now show that $K$ is perfect using the criterion of
Proposition \href{https://stacks.math.columbia.edu/tag/07LT}{proposition perfect is compact (Tag 07LT)}. Thus we let
$E_j \in D(R)$ be a family of objects parametrized by a set $J$.
We choose complexes $E_j^\bullet$ with flat terms
representing $E_j$, see for example Lemma \href{https://stacks.math.columbia.edu/tag/06Y4}{lemma K flat resolution (Tag 06Y4)}.
It is clear that
$$
0 \to
E_j^\bullet \to
E_j^\bullet/IE_j^\bullet \oplus E_j^\bullet/JE_j^\bullet \to
E_j^\bullet/(I + J)E_j^\bullet \to 0
$$
is a short exact sequence of complexes. Taking direct sums we obtain a
similar short exact sequence
$$
0 \to
\bigoplus E_j^\bullet \to
\bigoplus E_j^\bullet/IE_j^\bullet \oplus E_j^\bullet/JE_j^\bullet \to
\bigoplus E_j^\bullet/(I + J)E_j^\bullet \to 0
$$
(Note that $- \otimes_R R/I$ commutes with direct sums.)
This short exact sequence determines a distinguished triangle in $D(R)$, see
Derived Categories, Lemma \href{https://stacks.math.columbia.edu/tag/0152}{lemma derived canonical delta functor (Tag 0152)}.
Apply the homological functor $\Hom_{D(R)}(K, -)$ (see
Derived Categories, Lemma \href{https://stacks.math.columbia.edu/tag/0149}{lemma representable homological (Tag 0149)})
to get a commutative diagram
$$
\xymatrix{
\bigoplus \Hom_{D(R)}(K^\bullet, E_j^\bullet/(I + J))[-1] \ar[r] \ar[d] &
\Hom_{D(R)}(K^\bullet, \bigoplus E_j^\bullet/(I + J))[-1] \ar[d] \\
\bigoplus \Hom_{D(R)}(K^\bullet, E_j^\bullet/I \oplus E_j^\bullet/J)[-1]
\ar[r] \ar[d] &
\Hom_{D(R)}(K^\bullet, \bigoplus E_j^\bullet/I \oplus E_j^\bullet/J)[-1]
\ar[d] \\
\bigoplus \Hom_{D(R)}(K^\bullet, E_j^\bullet) \ar[r] \ar[d] &
\Hom_{D(R)}(K^\bullet, \bigoplus E_j^\bullet) \ar[d] \\
\bigoplus \Hom_{D(R)}(K^\bullet, E_j^\bullet/I \oplus E_j^\bullet/J)
\ar[r] \ar[d] &
\Hom_{D(R)}(K^\bullet, \bigoplus E_j^\bullet/I \oplus E_j^\bullet/J)
\ar[d] \\
\bigoplus \Hom_{D(R)}(K^\bullet, E_j^\bullet/(I + J)) \ar[r] &
\Hom_{D(R)}(K^\bullet, \bigoplus E_j^\bullet/(I + J))
}
$$
with exact columns. It is clear that, for any complex $E^\bullet$
of $R$-modules we have
\begin{align*}
\Hom_{D(R)}(K^\bullet, E^\bullet/I) & =
\Hom_{K(R)}(K^\bullet, E^\bullet/I) \\
& =
\Hom_{K(R/I)}(K^\bullet/IK^\bullet, E^\bullet/I) \\
& =
\Hom_{D(R/I)}(K^\bullet/IK^\bullet, E^\bullet/I)
\end{align*}
and similarly for when dividing by $J$ or $I + J$, see
Derived Categories,
Lemma \href{https://stacks.math.columbia.edu/tag/064B}{lemma morphisms from projective complex (Tag 064B)}.
Derived Categories. Thus all the horizontal
arrows, except for possibly the middle one, are isomorphisms as the complexes
$K^\bullet/IK^\bullet$, $K^\bullet/JK^\bullet$, $K^\bullet/(I + J)K^\bullet$
are perfect complexes of $R/I$, $R/J$, $R/(I + J)$-modules, see
Proposition \href{https://stacks.math.columbia.edu/tag/07LT}{proposition perfect is compact (Tag 07LT)}.
It follows from the $5$-lemma (Homology, Lemma \href{https://stacks.math.columbia.edu/tag/05QB}{lemma five lemma (Tag 05QB)})
that the middle map is an isomorphism and the lemma follows by
Proposition \href{https://stacks.math.columbia.edu/tag/07LT}{proposition perfect is compact (Tag 07LT)}.
\end{proof}
\end{document}
